% Laget automatisk av elevark/lag_elevark.py fra opplegg/repetisjon-bevegelsesmengde.html. Endringer her blir overskrevet.
\documentclass[11pt]{article}
\usepackage{elevark}
\begin{document}
\elevarkhode{Fysikk 1 · Bevegelsesmengde og støt}{Repetisjon: bevegelsesmengde og støt}{https://lektorbo.no/opplegg/repetisjon-bevegelsesmengde.html}{Opplegget på nett}

I et støt er kreftene store og kortvarige, og vanskelige å måle. Men bevegelsesmengden er bevart, og det gjør støt enkle å regne på. Her er en oversikt og blandede oppgaver med løsningsforslag.

\section*{Tre typer støt}

\begin{center}\small
\begin{tabularx}{\linewidth}{YYY}
\toprule
\textbf{Støt} & \textbf{Bevegelsesmengde} & \textbf{Kinetisk energi} \\
\midrule
Elastisk & bevart & bevart \\
Uelastisk & bevart & noe blir til varme, lyd og deformasjon \\
Fullstendig uelastisk (henger sammen etterpå) & bevart & størst mulig tap \\
\bottomrule
\end{tabularx}
\end{center}

Bevegelsesmengden er bevart når de ytre kreftene er små i forhold til kreftene i støtet. Husk at bevegelsesmengde er en vektor. Velg en positiv retning, og gi farten minus når den peker motsatt vei.

\section*{Formler}

\begin{itemize}
\item Bevegelsesmengde: $p = mv$.
\item Impuls: $I = FΔt = Δp$. Med varierende kraft er impulsen arealet under F,t-grafen.
\item Bevaring: $m_{1}v_{1} + m_{2}v_{2} = m_{1}u_{1} + m_{2}u_{2}$.
\item Kinetisk energi: $E_{\mathrm{k}} = \tfrac{1}{2}mv^{2} = p^{2}/(2m)$.
\end{itemize}

\section*{Oppgaver}

Bruk $g = 9{,}81\,\mathrm{m/s^{2}}$. Velg positiv retning i den første fartsretningen.

\begin{oppgaveliste}
\oppgave \fase{forklar}{Begrep}Er påstandene sanne eller usanne? Begrunn svaret.

a) Bevegelsesmengden er bare bevart i elastiske støt.\newline b) En ball som spretter rett tilbake fra en vegg, får større endring i bevegelsesmengde enn en ball som stopper.\newline c) En kollisjonspute gjør at impulsen på hodet blir mindre.\newline d) Når et gevær skytes, får geværet og kula like stor bevegelsesmengde.\newline e) $I$ et fullstendig uelastisk støt forsvinner all kinetisk energi.
\svar
\oppgave \fase{test}{Regn}En bil på 1500~kg i 20~m/s kjører inn i en bil på 1000~kg som står i ro. Bilene henger sammen etterpå. Finn farten rett etter støtet og hvor stor del av den kinetiske energien som går tapt. Sjekk med simuleringen.\footnote{\url{lektorbo.no/sim/vognstot.html}}
\svar
\oppgave \fase{test}{Regn}En fotball på 0,43~kg ligger i ro og blir sparket. Den får farten 25~m/s, og foten er i kontakt med ballen i 8,0~ms. Finn impulsen og den gjennomsnittlige kraften. Se Myk landing.\footnote{\url{lektorbo.no/sim/impuls.html}}
\svar
\oppgave \fase{test}{Regn}Vogn A på 0,50~kg med farten 0,60~m/s støter elastisk mot vogn B på 0,25~kg som står i ro. Etter støtet har A farten 0,20~m/s i samme retning. Finn farten til B, og vis at den kinetiske energien er bevart.
\svar
\oppgave \fase{test}{Regn}En kanon på 800~kg står fritt og skyter en kule på 5,0~kg med farten 300~m/s. Hvilken fart får kanonen? Hvor mye kinetisk energi får kula og kanonen? Sjekk med simuleringen.\footnote{\url{lektorbo.no/sim/rekyl.html}}
\svar
\oppgave \fase{test}{Regn}En kule på 10~g blir skutt inn i en trekloss på 2,0~kg som henger i snorer. Klossen og kula svinger opp 5,0~cm. Finn farten til klossen rett etter støtet og farten til kula. Hvorfor kan vi ikke bruke energibevaring for selve støtet? Sjekk med simuleringen.\footnote{\url{lektorbo.no/sim/ballistisk-pendel.html}}
\svar
\oppgave \fase{utfordring}{Sammensatt}En ball på 0,15~kg ligger i ro og blir slått med et balltre. Kraften stiger jevnt fra 0 til 600~N og faller jevnt til 0 igjen. Hele slaget varer 0,020~s. Skisser F,t-grafen. Finn impulsen og farten til ballen. Hvor stor er den gjennomsnittlige kraften?
\svar
\end{oppgaveliste}

\end{document}
